(1) 动能 \(T=\frac12I_1(\dot\theta^2+\dot\phi^2\sin^2\theta)+\frac12I_3(\dot\psi+\dot\phi\cos\theta)^2\)；由于质心高度 \(z=\ell\cos\theta\)，重力势能 \(U=Mg\ell\cos\theta\)，所以拉氏量为 \(L=T-U\)，固定点反力不对轴做功。
(2) 在规则进动中 \(\dot\theta=\ddot\theta=0\)。对 \(\theta\) 的 Euler–Lagrange 方程为 \(\partial L/\partial\theta=0\)，即 \(I_1\omega_p^2\sin\theta\cos\theta-I_3s\omega_p\sin\theta+Mg\ell\sin\theta=0\)。除以 \(\sin\theta>0\) 得 \(I_1\cos\theta\,\omega_p^2-I_3s\,\omega_p+Mg\ell=0\)。
(3) 二次方程的两根 \(\omega_{p,\pm}=[I_3s\pm\sqrt{I_3^2s^2-4I_1Mg\ell\cos\theta}]/(2I_1\cos\theta)\)，在 \(s>0\) 下分别对应快、慢进动。存在实规则进动要求 \(I_3^2s^2\ge4I_1Mg\ell\cos\theta\)；若不满足，指定倾角无法维持恒定进动而出现章动。
(4) 在 \(I_3^2s^2\gg4I_1Mg\ell\cos\theta\) 时，对较小根作平方根展开，得 \(\omega_{p,-}\simeq Mg\ell/(I_3s)\)，而快根约为 \(I_3s/(I_1\cos\theta)\)。\([Mg\ell/(I_3s)]=(\mathrm{kg\,m^2\,s^{-2}})/(\mathrm{kg\,m^2\,s^{-1}})=\mathrm{s^{-1}}\)，量纲正确。
